\section{Алгебраическая и геометрическая кратности собственного числа. Теорема о связи алгебраической и геометрической кратностей}

\begin{definition}
Пусть $\lambda_0$ --- собственное число линейного оператора $\varphi$.
\begin{enumerate}
    \item \textbf{Алгебраической кратностью} собственного числа $\lambda_0$ (обозначается $m_a$ или $m$) называется его кратность как корня характеристического многочлена $P_\varphi(\lambda) = \det(A_\varphi - \lambda E)$.
    \item \textbf{Геометрической кратностью} собственного числа $\lambda_0$ (обозначается $m_g$ или $k$) называется размерность соответствующего собственного подпространства: $m_g = \dim L_{\lambda_0}$.
\end{enumerate}
\end{definition}

\begin{theorem}[О связи алгебраической и геометрической кратностей]
Для любого собственного числа $\lambda_0$ линейного оператора $\varphi$ его геометрическая кратность не превосходит алгебраической:
\[
m_g \le m_a.
\]
\end{theorem}

\begin{proof}
Пусть $m_g = \dim L_{\lambda_0} = k$. Выберем в подпространстве $L_{\lambda_0}$ базис $e_1, e_2, \dots, e_k$. 
По теореме о дополнении базиса подпространства до базиса всего пространства, мы можем добавить векторы $e_{k+1}, \dots, e_n$ так, чтобы система $E = \{e_1, \dots, e_k, e_{k+1}, \dots, e_n\}$ образовала базис всего пространства $V$ (где $n = \dim V$).

Найдем матрицу оператора $\varphi$ в этом базисе $E$. 
Для первых $k$ векторов, так как они принадлежат $L_{\lambda_0}$, имеем:
\[
\varphi(e_i) = \lambda_0 e_i, \quad i = 1, \dots, k.
\]
Для остальных векторов $\varphi(e_j)$ ($j > k$) раскладываются по всему базису $E$. Следовательно, матрица $A_\varphi$ в базисе $E$ имеет блочную структуру:
\[
A_\varphi = \begin{pmatrix}
\lambda_0 E_k & B \\
0 & C
\end{pmatrix},
\]
где $E_k$ --- единичная матрица порядка $k$, $B$ --- некоторая матрица размера $k \times (n-k)$, $0$ --- нулевая матрица размера $(n-k) \times k$, а $C$ --- квадратная матрица порядка $n-k$.

Вычислим характеристический многочлен оператора $\varphi$:
\[
P_\varphi(\lambda) = \det(A_\varphi - \lambda E_n) = \det \begin{pmatrix}
(\lambda_0 - \lambda) E_k & B \\
0 & C - \lambda E_{n-k}
\end{pmatrix}.
\]
Используя свойство определителя блочной треугольной матрицы, получаем:
\[
P_\varphi(\lambda) = \det((\lambda_0 - \lambda) E_k) \cdot \det(C - \lambda E_{n-k}) = (\lambda_0 - \lambda)^k \cdot \det(C - \lambda E_{n-k}).
\]
Из этого разложения видно, что многочлен $P_\varphi(\lambda)$ делится на $(\lambda_0 - \lambda)^k$. Следовательно, $\lambda_0$ является корнем характеристического уравнения кратности не менее $k$. То есть алгебраическая кратность $m_a \ge k = m_g$.
\hfill$\blacksquare$
\end{proof}
